\documentclass{article}
\usepackage{apacite}
\usepackage{hyperref}
\usepackage{multicol}
\usepackage{natbib}
\usepackage{sophia}
\usepackage[utf8]{inputenc}
\usepackage{array}
\setlength{\parindent}{0pt}
\usepackage{hyperref}

\usetikzlibrary{calc,patterns,angles,quotes}
\usepackage{pgfplots}
\pgfplotsset{width=15cm,height=10cm,compat=1.9}

\usepackage[english]{babel}
\usepackage{fancyhdr}
\usepackage{lastpage}
\pagestyle{fancy}
\fancyhf{}


\pagestyle{empty}% for cropping
\usepackage{tikz}
\usetikzlibrary{calc}
\usepackage{zref-savepos}
\usepackage{tabu}

\newcounter{NoTableEntry}
\renewcommand*{\theNoTableEntry}{NTE-\the\value{NoTableEntry}}

\newcommand*{\strike}[2]{%
  \multicolumn{1}{#1}{%
    \stepcounter{NoTableEntry}%
    \vadjust pre{\zsavepos{\theNoTableEntry t}}% top
    \vadjust{\zsavepos{\theNoTableEntry b}}% bottom
    \zsavepos{\theNoTableEntry l}% left
    \hspace{0pt plus 1filll}%
    #2% content
    \hspace{0pt plus 1filll}%
    \zsavepos{\theNoTableEntry r}% right
    \tikz[overlay]{%
      \draw
        let
          \n{llx}={\zposx{\theNoTableEntry l}sp-\zposx{\theNoTableEntry r}sp-\tabcolsep},
          \n{urx}={\tabcolsep},
          \n{lly}={\zposy{\theNoTableEntry b}sp-\zposy{\theNoTableEntry r}sp},
          \n{ury}={\zposy{\theNoTableEntry t}sp-\zposy{\theNoTableEntry r}sp}
        in
        (\n{llx}, \n{lly}) -- (\n{urx}, \n{ury})
      ;
    }% 
  }%
}

\usepackage{tikz}
\usetikzlibrary{calc}

\tikzset{
    right angle quadrant/.code={
        \pgfmathsetmacro\quadranta{{1,1,-1,-1}[#1-1]}     % Arrays for selecting quadrant
        \pgfmathsetmacro\quadrantb{{1,-1,-1,1}[#1-1]}},
    right angle quadrant=1, % Make sure it is set, even if not called explicitly
    right angle length/.code={\def\rightanglelength{#1}},   % Length of symbol
    right angle length=2ex, % Make sure it is set...
    right angle symbol/.style n args={3}{
        insert path={
            let \p0 = ($(#1)!(#3)!(#2)$) in     % Intersection
                let \p1 = ($(\p0)!\quadranta*\rightanglelength!(#3)$), % Point on base line
                \p2 = ($(\p0)!\quadrantb*\rightanglelength!(#2)$) in % Point on perpendicular line
                let \p3 = ($(\p1)+(\p2)-(\p0)$) in  % Corner point of symbol
            (\p1) -- (\p3) -- (\p2)
        }
    }
}
\pagestyle{fancy}
\fancyhf{}
\title{Andi's Solution To Fiona's Advanced Function Test 1}
\author{Andi Xie}
\date{\today}
\lhead{Andi Xie}
\rhead{\today}
\pagestyle{fancy}

 
\cfoot{Page \thepage \hspace{1pt} of \pageref{LastPage}}


\begin{document}

\maketitle
\tableofcontents
\newpage


\section{Question 1}
\begin{custom-big}[Part a]
\begin{align*}
    \lim_{x\to -\infty}f(x)&=\infty\\
    \lim_{x\to \infty}f(x)&=-\infty
\end{align*}
\end{custom-big}




\begin{custom-big}[Part b]
\[
-6x^3+8x^2+10x-17=(2x^2-3)(-3x+4)+(x-5)
\]
Andi chose to not type out the long division because he's lazy (he doesn't know how to)
\end{custom-big}




\section{Question 2}
\begin{custom-big}[Part a]
\begin{align*}
5x^3-12x^2+x+6&=0\\
(x-1)(x-2)(5x+3)&=0\\
x&=-\frac 35,1,2
\end{align*}
\end{custom-big}

\begin{custom-big}[Part b]
\begin{align*}
(x^3+1)(x^3-8)&=0\\
x^3&=-1,8\\
x&=-1,2,\frac{1\pm\sqrt{3}i}{2},-1\pm\sqrt{3}i
\end{align*}
\end{custom-big}


\section{Question 3}
Andi doesn't know, cryyyyy ;-;


\section{Question 4}
\begin{custom-big}[Part a]
\begin{align*}
    -2x^3+3x^2+5x-6&\ge 0\\
    (x-1)(2x+3)(-x+2)&\ge 0\\
    x\leq -\frac 32,1\leq x\leq 2
\end{align*}
\end{custom-big}

\begin{custom-big}[Part b]
When $x<2$ or $x>4$:\\
\begin{align*}
    (3x+1)(x-2)&<(x+1)(x-4)\\
    x^2-x+1&<0
\end{align*}
No Solution.\\
When $2<x<4$:\\
\begin{align*}
    x^2-x+1&>0
\end{align*}
This is always true, therefore $2<x<4$
\end{custom-big}


\section{Question 5}
\begin{custom-big}[Solution]
From the remainder theorem:
\[
\begin{cases}
f(-2)=-49\\
f(1)=-1
\end{cases}
\]
Solving the equation gives $a=-1,b=3$
\end{custom-big}

\pagebreak
\section{Question 6}
\begin{custom-big}[Solution]
\[
f(x)=\frac{-2x^3(x-2)}{x\left(x-\frac 34\right)\left(x+\frac 23\right)(x-2)}\]
\end{custom-big}

\section{Question 7}
\begin{custom-big}[Solution]
\[
f(x)=-6(x-2)(x+3)(x+1)^2\]
\end{custom-big}



\section{Question 6}
\begin{custom-big}[Part a]
\[
y=-4\left(x+1\right)\left(x-\frac 12\right)\left(x-2\right)^3\]
\end{custom-big}

\begin{custom-big}[Part b]
\[
y=4\left(\frac 12 x+1\right)\left(\frac 12 x-\frac 12\right)\left(\frac 12 x-2\right)^3\]
\end{custom-big}

\begin{custom-big}[Part c]
Horizontal expansion by factor of 2, flip over $x$-axis
\end{custom-big}

\pagebreak
\section{Question 9}
\begin{custom-big}[Part a]
\[\frac{(x-2)(x+2)(x+3)}{(x-2)(x-3)}\]
No HA
VA at $x=3$\\
Hole at $(2,-20)$
\end{custom-big}

\begin{custom-big}[Part b]
Desmos ftw
\end{custom-big}


\section{Question 10}
\begin{custom-big}[Solution]
\[\frac{(x-2)^2}{x(x-3)}\]
HA at $y=1$\\
VA at $x=0,x=3$\\
$x$-interception at $x=2$\\
No holes
\end{custom-big}


\section{Question 11}
\begin{custom-big}[Part a]
\[f(x)=-(x-4)(x+3)^3\]
\end{custom-big}

\begin{custom-big}[Part b]
\[f(x)=\left(x+1\right)\left(x-1\right)^{2}\left(x^{2}-2x\right)+0.679552342168\]
PS. Andi cheated by using Calculus
\end{custom-big}


\section{Bonus}
\begin{custom-big}[Solution]
\begin{align*}
    \left(x+\frac 1x\right)^3=x^3+3x+3\frac 1x+\frac{1}{x^3}&\implies x^3+\frac{1}{x^3}=\left(x+\frac 1x\right)^3-3\left(x+\frac 1x\right)\\
    \left(x+\frac 1x\right)^2=x^2+2+\frac{1}{x^2}&\implies x^2+\frac{1}{x^2}=\left(x+\frac 1x\right)^2-2
\end{align*}
Let $x+\dfrac 1x=y$:
\begin{align*}
    y^3-3y+y^2-2+y&=28\\
    y^3+y^2-2y-30&=0\\
    (y-3)(y^2+4y+10)&=0\\
    \therefore y&=3\\
    x+\frac 1x&=3\\
    x^2-3x+1&=0\\
    x&=\frac{3\pm\sqrt 5}{2}
\end{align*}
\end{custom-big}
\end{document}
